\documentclass[aps,prl,twocolumn,superscriptaddress,floatfix]{revtex4-2}
\usepackage{amsmath,amssymb,amsfonts,booktabs,array,xcolor,hyperref}

\begin{document}

\title{BCT Letter 84:\\
The BCT Periodic Table: A Topological Reclassification of the Elements}

\author{Michel Robert Cabri\'{e}}
\affiliation{Independent Researcher, Barrys Reef, Victoria, Australia}
\affiliation{ORCID: 0009-0007-9561-9859}
\date{20 March 2026}

\begin{abstract}
The standard periodic table reveals periodicity but obscures its origin.
The BCT Superfluid Lattice Model provides a complete first-principles
derivation of the table's structure from a single geometric postulate.
In this framework, each element's chemical character is a direct consequence
of its Hopf multiplet topology --- specifically the fill fraction of its
Octet-Hopfion Condensate (OHC) modes. We present two novel arrangements
of all 118 elements organised by this topology, and identify three
\textit{topological axes} invisible in any prior periodic table:
the Life Axis ($S_H = 1$ half-fill elements, H$\to$C$\to$Si$\to$Ge$\to$Sn$\to$Pb$\to$Fl),
the Noble Metal Axis (d-shell closure attractors, Cu/Ag/Au), and the
d/f Half-Fill Axis (Hopf symmetry points, Cr/Mo/Re and Eu/Am).
These axes are derived consequences of the OHC mode structure proven
in Letters 78--83, not empirical observations.
Companion appendices: KC1 (full classification), KC2 (axis proofs).
\end{abstract}

\maketitle

\section{Introduction}

Mendel\'eev's 1869 arrangement of elements was a triumph of pattern
recognition. But it was descriptive, not explanatory. The question
\textit{why} elements repeat in periods of 2, 8, 8, 18, 18, 32, 32
was not answered by the table itself; it required quantum mechanics
and the Pauli exclusion principle, which in turn required explanation.

The BCT Superfluid Lattice Model resolves this in a single stroke.
The quantum vacuum is a BCT superfluid lattice with $c/a = \sqrt{2}$.
The OHC is its ground state. Electron shells are OHC Bessel resonance
modes. The shell capacity $2n^2$ is a derived theorem (Letter~81).
Pauli exclusion is a derived theorem (Letter~81). Noble gas inertness
is a derived theorem (Letter~81). The anomalous configurations of
Cr, Mo, Cu, Ag, Au, Eu, and Yb are all consequences of Hopf symmetry
attractors (Letters~82--83). The uniqueness of carbon as the basis
of life follows from $S_H = 1$ at $n=2$ (Letter~82).

\section{BCT Element Classification}

Each element is assigned to one of eleven BCT topological classes
(see Appendix KC1 for full classification):

\begin{center}
\begin{tabular}{p{22mm}p{50mm}}
\toprule
Class & Defining BCT property \\
\midrule
Noble gas & Closed Hopf multiplet. Bonding impossible (L81). \\
Life axis & $S_H = 1$ at sp$^3$ half-fill. Unique max (L82). \\
d half-fill & $S_H(d) = 0.50$. Hopf attractor (L83). \\
d closure & Closed d multiplet $N_d = 10$ (L83). \\
f half-fill & $S_H(f) = 1.00$. f Hopf attractor (KB3). \\
f closure & Closed f multiplet $N_f = 14$ (KB3). \\
Ferromagnetic & Collective Hopf alignment (L83). \\
s-block & Outermost mode is $r_\mathrm{oct}$. \\
p-block & Outermost modes are $r_\mathrm{tet}$. \\
d-block & Active d-modes, not at symmetry point. \\
f-block & Active f-modes. \\
\bottomrule
\end{tabular}
\end{center}

\section{The Janet Left-Step Layout}

The Janet (left-step) layout reads left-to-right as
f $\to$ d $\to$ p $\to$ s, corresponding exactly to the OHC Bessel
resonance filling sequence from deepest mode to outermost.
This is the natural BCT layout. Helium sits above Beryllium
(both $s^2$ closures), not above Neon. The noble gas column
forms the right edge of every period --- marking topological
closure of each Hopf shell.

\section{The Circular Layout and Three Topological Axes}

When elements are arranged on concentric rings (one per Hopf
shell, $n=1$ inner to $n=7$ outer) with noble gases at
12~o'clock, three topological axes become visible as geometric
features:

\subsection{The Life Axis ($\theta = 180^\circ$, 6~o'clock)}

All elements with $S_H = 1$ at sp$^3$ half-fill appear at
$\theta = 180^\circ$ in their ring.
For periods~1--3 this is exact:
\begin{equation}
  \theta_C = \frac{4}{8} \times 360^\circ = 180^\circ
\end{equation}
The Life Axis: H$\to$C$\to$Si$\to$Ge$\to$Sn$\to$Pb$\to$Fl.
Carbon's uniqueness as the basis of known life is a consequence
of occupying this axis at $n=2$, the lowest available shell
for organic chemistry.

\subsection{The Noble Metal Axis ($\theta \approx 220^\circ$)}

All d-shell closure elements (Cu, Ag, Au, Cn) appear at
$\theta = (11/18) \times 360^\circ \approx 220^\circ$
in their rings. This is period-independent, fixed by the
d-block capacity and closure attractor position.

\subsection{The d/f Half-Fill Axis ($\theta \approx 140^\circ$)}

All d-half-fill elements (Cr, Mo, Re) appear at
$\theta = (7/18) \times 360^\circ \approx 140^\circ$.
The f-half-fill elements (Eu, Am) occupy the analogous
position in the f-block sub-sequence.

\section{Key Predictions}

\begin{center}
\begin{tabular}{lll}
\toprule
Prediction & BCT & Obs./status \\
\midrule
H$_2$O angle & $104.478^\circ$ & $104.4776^\circ$ (L80) \\
Carbon $S_H$ & 1.000 (unique) & Confirmed (L82) \\
Cu config & $3d^{10}4s^1$ & $3d^{10}4s^1$ (L83) \\
Au 521\,nm & $5d{\to}6s$ & 521\,nm (L83) \\
Hg liquid & Double Hopf closure & Liquid RT (L83) \\
\bottomrule
\end{tabular}
\end{center}

\section{Conclusion}

The BCT Periodic Table is the first arrangement in which every
structural feature --- period lengths, noble gas positions,
anomalous configurations, gold's colour --- has a derivation
rather than a description. The three topological axes are not
imposed from outside; they emerge geometrically when elements
are placed by Hopf topology. Life found the only available
topological axis. The BCT programme explains why.

Companion figures available on Zenodo:
\texttt{BCT\_PeriodicTable.pdf} (Janet layout) and
\texttt{BCT\_PeriodicTable\_Round\_v12.pdf} (circular layout).

\begin{thebibliography}{99}
\bibitem{L78} M.~R.~Cabri\'{e}, BCT Letter~78, Zenodo (2026).
\bibitem{L80} M.~R.~Cabri\'{e}, BCT Letter~80, Zenodo (2026).
\bibitem{L81} M.~R.~Cabri\'{e}, BCT Letter~81, Zenodo (2026).
\bibitem{L82} M.~R.~Cabri\'{e}, BCT Letter~82, Zenodo (2026).
\bibitem{L83} M.~R.~Cabri\'{e}, BCT Letter~83, Zenodo (2026).
\bibitem{KC1} M.~R.~Cabri\'{e}, BCT Appendix KC1, Zenodo (2026).
\bibitem{KC2} M.~R.~Cabri\'{e}, BCT Appendix KC2, Zenodo (2026).
\end{thebibliography}

\end{document}
